\documentclass{chsmc}
\chsmcsetup{
	year = 2013,
	part = 1,
	type = B,
	show-answers = false,
}
\begin{document}

\section{填空题：本大题共 8 小题，每小题 8 分，共 64 分.}

\begin{enumerate}
	\item 已知锐角三角形的三条边长都是整数，其中两条边长分别为 $3$ 和 $4$，
		则第三条边的边长为 \blankline
	\item 设 $\ii = \sqrt{-1}$ 为虚数单位，则
		$\ii + 2\ii^2 + 3\ii^3 + \cdots + 2013\ii^{2013} =$ \blankline
	\item 设集合 $A = \{2,0,1,3\}$，集合
		$B = \{x\mid -x \in A,\, 2 - x^2 \notin A\}$.
		则集合 $B$ 中所有元素的和为 \blankline
	\item 已知正三棱锥 $P$-$ABC$ 底面边长为 $1$，高为 $\sqrt{2}$，
		则其内切球半径为 \blankline
	\item 在区间 $[0,\pi)$ 中，方程 $\sin 12x = x$ 解的个数为 \blankline
	\item 定义在 $\mathbf{R}$ 上的函数
		$f(x) = \dfrac{\sin \pi x}{\sqrt{1+x+x^2}}$ 的最小值是 \blankline
	\item 设 $a$, $b$ 为实数，函数 $f(x) = ax + b$ 满足：
		对任意 $x \in [0,1]$，有 $|f(x)| \leqslant 1$.
		则 $ab$ 的最大值为 \blankline
	\item 将正九边形的每个顶点等概率地涂上红、蓝两种颜色之一，
		则存在三个同色的顶点构成锐角三角形的概率为 \blankline
\end{enumerate}

\section{解答题：本大题共 3 个小题，满分 56 分. 解答应写出文字说明、证明过程或演算步骤.}

\begin{enumerate}[resume]
	\item (本题满分 16 分) 已知数列 $\{a_n\}$ 满足：
		$a_1 = 2$，$a_n = 2(n+a_{n-1})$，$n = 2,3,\cdots$，
		求数列 $\{a_n\}$ 的通项公式.
	\item (本题满分 20 分) 假设 $a$, $b$, $c > 0$，且 $abc = 1$，
		证明：$a+b+c \leqslant a^2+b^2+c^2$.
	\item (本题满分 20 分) 在平面直角坐标系 $xOy$ 内，点 $F$ 的坐标为 $(1,0)$，
		点 $A$, $B$ 在抛物线 $y^2 = 4x$ 上，满足
		$\overrightarrow{OA}\cdot\overrightarrow{OB} = -4$，
		$\bigl|\overrightarrow{FA}\bigr| - \bigl|\overrightarrow{FB}\bigr| = 4\sqrt{3}$，
		求 $\overrightarrow{FA}\cdot\overrightarrow{FB}$ 的值.
\end{enumerate}

\end{document}
